\section*{Bab \ref{ch:b2:cell-signalling} --- Kurir Kimia dan Transduksi Isyarat}

\begin{solution}{exo:b2:cell-signalling:1}
Insulin: endokrin, larut air (sebuah peptida). Kortisol: endokrin,
larut lipid. Asetilkolin pada sebuah sinapsis: sinaptik, larut air.
Histamin pada sebuah luka: parakrin, larut air. Estradiol: endokrin,
larut lipid. Adrenalin dari adrenalnya: endokrin, larut air. Nitrogen
monoksida: parakrin, larut lipid (sebuah gas yang menyeberangi
membran).
\end{solution}

\begin{solution}{exo:b2:cell-signalling:2}
Adrenalin mengikat reseptornya (padam: disosiasi, penginternalan
reseptornya); reseptornya mengaktifkan \omterm{prop:b2:cell-signalling:gpcr}{protein G}, GDP menjadi GTP
(padam: hidrolisis GTP dalam hitungan detik); \omterm{prop:b2:cell-signalling:gpcr}{protein G-nya}
mengaktifkan adenilil siklase, yang membuat cAMP (padam:
fosfodiesterase); cAMP mengaktifkan kinase A (padam: kehilangan cAMP);
kinase A memfosforilasi fosforilase kinase, yang memfosforilasi
glikogen fosforilase (padam: fosfatase); fosforilase melepaskan
glukosa-1-fosfat dari glikogennya, lalu diubah dan diekspor sebagai
glukosa.
\end{solution}

\begin{solution}{exo:b2:cell-signalling:3}
Reseptor permukaan mengikat kurir yang larut air pada membrannya lalu
bekerja dalam hitungan detik lewat \omterm{prop:b2:cell-signalling:gpcr}{kurir kedua} dan kinase, sehingga
mengubah keaktifan protein yang sudah ada. \omterm{prop:b2:cell-signalling:nuclear}{Reseptor inti} mengikat kurir
yang larut lipid di dalam selnya lalu bekerja dalam hitungan jam
sebagai \omterm{prop:b2:cell-differentiation:transcriptional}{faktor transkripsi}, sehingga mengubah protein mana yang dibuat.
\end{solution}

\begin{solution}{exo:b2:cell-signalling:4}
\omterm{prop:b2:cell-signalling:gpcr}{AMP siklik}: dibuat adenilil siklase, dimusnahkan fosfodiesterase.
Kalsium: masuk lewat kanal atau meninggalkan retikulum endoplasmanya
sebagai jawaban atas inositol trifosfat, disingkirkan oleh pompa.
Inositol trifosfat (bersama diasilgliserol): dibuat fosfolipase C dari
sebuah lipid membran, disingkirkan oleh fosfatase.
\end{solution}

\begin{solution}{exo:b2:cell-signalling:5}
$\theta = [L]/(K_d + [L])$: $0.09$, $0.5$, $0.91$, $0.99$.
Melipatduakan reseptornya membiarkan bagiannya tak berubah dan
melipatduakan jumlah yang terisi --- dan dengan demikian tanggapannya.
\end{solution}

\begin{solution}{exo:b2:cell-signalling:6}
$50\times 200\times 100\times 300 = 3\times 10^{8}$. Seratus reseptor
yang terisi memberikan $3\times 10^{10}$ molekul hasil tiap detik ---
sebuah batas atas, karena tahapnya menjenuh.
\end{solution}

\begin{solution}{exo:b2:cell-signalling:7}
\omterm{prop:b2:cell-signalling:gpcr}{Protein G} yang terkunci itu menjaga siklasenya tetap menyala, cAMP-nya
tetap tinggi, kinase A menjaga kanal klorida sel ususnya tetap
terfosforilasi dan terbuka, kloridanya disekresikan, lalu natrium dan
air menyusulnya ke dalam lumennya --- berliter-liter tiap hari.
Kuncinya adalah sebuah perubahan kovalen pada \omterm{prop:b2:cell-signalling:gpcr}{protein G-nya}, yang hanya
dibatalkan ketika selnya sendiri diganti, berhari-hari kemudian.
\end{solution}

\begin{solution}{exo:b2:cell-signalling:8}
$0.9\times 10^{-6}\,\text{mol/L}\times 2\times 10^{-12}\,\text{L} =
1.8\times 10^{-18}$ mol, yaitu $1.1\times 10^{6}$ ion (sebenarnya lebih
banyak, karena penyangganya mengikat sebagian besar yang masuk). Satu
kanal pada $10^{6}$ tiap detik akan memerlukan sedetik; seribu kanal
mengerjakannya dalam semilidetik, yaitu skala waktu yang diperlukan
sebuah isyarat.
\end{solution}

\begin{solution}{exo:b2:cell-signalling:9}
Kedua lintasannya bertemu pada enzim yang sama: kinase A glukagonnya
memfosforilasi glikogen sintase (padam) dan fosforilase kinase
(nyala); riam insulinnya mengaktifkan fosfatase yang menyingkirkan
fosfatnya. Keadaan enzimnya ditetapkan oleh perimbangan keaktifan
kinase dan fosfatasenya, sehingga yang berarti adalah nisbah kedua
\omterm{def:b2:cell-signalling:kinds}{hormonnya}; ketika keduanya tinggi fosfatasenya biasanya menang dan
glikogennya disimpan.
\end{solution}

\begin{solution}{exo:b2:cell-signalling:10}
Adrenalin: cAMP berdifusi melintasi sebuah sel \qty{20}{\micro m} dalam
kurang dari sedetik ($x^{2}/2D$ dengan $D \approx \qty{3e-10}{m^2/s}$),
dan tiap enzim riamnya sudah ada --- beberapa detik. Kortisol: sebuah
transkrip memerlukan \qty{60}{s}, sebuah protein \qty{100}{s}; dari
beberapa lusin mRNA yang masing-masing membawa selusin ribosom, selnya
membuat beberapa ratus molekul enzim tiap menit, sehingga beribu yang
diperlukan bagi sebuah efek yang terukur memerlukan berjam-jam, setelah
beberapa menit yang diperlukan kortisol untuk masuk dan mencapai DNA-nya.
\end{solution}

\begin{solution}{exo:b2:cell-signalling:11}
Reseptornya memberi isyarat terus-menerus: riam MAP kinasenya
mendorong selnya untuk membelah tanpa faktor pertumbuhan apa pun, dan
ia mengabaikan ketiadaan isyarat yang biasanya menahannya. Karena
sebuah \omterm{def:b2:genome-diversification:mutation}{mutasi} tunggal menyingkirkan sebuah kendali atas pembelahannya,
reseptor semacam itu (dan kinase di hilirnya) termasuk onkogen yang
paling lazim. Sebuah obat yang menempati tapak ATP kinasenya
menghentikan fosforilasinya lalu membungkam reseptornya --- yaitu asas
beberapa obat kanker.
\end{solution}

\begin{solution}{exo:b2:cell-signalling:12}
Endokrin: mencapai tiap sel dalam semenit, murah tiap selnya, tidak
dapat mengalamati satu sel, bekerja selama beberapa menit sampai
beberapa hari --- tak tergantikan bagi keadaan seluruh tubuhnya
(berpuasa, tumbuh, berbiak). Saraf: semilidetik, satu sasaran,
pengabelan yang mahal, bekerja selama beberapa milidetik --- tak
tergantikan bagi gerakan dan refleks yang cepat. Keduanya dipakai bagi
pesan yang sama ketika kecepatan dan jangkauan sama-sama diinginkan:
saraf simpatis dan medula adrenalnya sama-sama mengantarkan
noradrenalin atau adrenalin ketika ketakutan.
\end{solution}

\begin{solution}{pb:b2:cell-signalling:1}
\textbf{1.} $1\,\text{nmol}/5\,\text{L} = \qty{0.2}{nmol/L}$.
\textbf{2.} $\theta = 0.2/1.2 = 0.17$: kira-kira 3300 reseptor terisi.
\textbf{3.} Tiga paruh waktu: \qty{0.025}{nmol/L}; $\theta = 0.024$.
\textbf{4.} Kira-kira setengah menit sampai semenit, yaitu waktu yang
diperlukan \omterm{def:b2:blood-circulation:blood}{darahnya} untuk berkeliling; sarafnya mengantarkan
penghantarnya langsung ke hatinya dalam sedetik.
\textbf{5.} Bagian yang sama, tetapi $33\,000$ reseptor terisi: sebuah
tanggapan sepuluh kali lebih besar.
\textbf{6.} Pada \qty{1}{nmol/L} sebuah reseptor dengan $K_d = \qty{1}{\micro
mol/L}$ terisi seperseribu: tanpa tanggapan. Menyerasikan $K_d$ dengan
jangkauan yang beredar menaruh bagian kurva keterisian yang curam tepat
pada tempat \omterm{def:b2:cell-signalling:kinds}{hormonnya} sungguh berubah.
\textbf{7.} $20\times 100\times 1\times 50\times 50\times 1000 =
5\times 10^{9}$ molekul glukosa tiap detik tiap reseptor yang terisi.
\textbf{8.} $3300\times 20\times 100 = 6.6\times 10^{6}$ cAMP pada
detik pertama; di dalam \qty{5e-12}{L}, $1.1\times 10^{-17}$ mol,
kira-kira \qty{2}{\micro mol/L}.
\textbf{9.} $3300\times 5\times 10^{9} = 1.7\times 10^{13}$ tiap detik, $10^{15}$
tiap menit tiap selnya.
\textbf{10.} $2\times 10^{11}$ sel $\times 10^{15} = 2\times
10^{26}$ molekul tiap menit, 330 mol, \qty{6e4}{g} tiap menit ---
sepuluh ribu kali \qty{5}{g/min} yang teramati. Batinya hanya
berlipat selama tiap tahapnya belum jenuh; sesungguhnya fosforilasenya
dibatasi substratnya dan fosfatasenya, dan ekspor glukosanya dibatasi
pengangkutnya. Bati riamnya dibelanjakan untuk kecepatan dan kepekaan,
bukan untuk keluaran.
\textbf{11.} Pada \qty{5}{g/min}, dua puluh menit; ketakutannya usai
dalam hitungan menit dan sakelar pemadamnya menghentikan riamnya jauh
sebelum itu.
\textbf{12.} Hidrolisis GTP oleh \omterm{prop:b2:cell-signalling:gpcr}{protein G-nya} (detik);
fosfodiesterase pada cAMP-nya (detik); fosfatase pada enzim yang
terfosforilasi (detik sampai menit); pemanaan reseptornya (menit). Yang
tercepat adalah jam \omterm{prop:b2:cell-signalling:gpcr}{protein G-nya} sendiri dan fosfodiesterasenya.
\textbf{13.} cAMP-nya tidak dimusnahkan: ia menimbun dan bertahan,
kinasenya tetap menyala, dan keluaran glukosanya lebih besar serta
bertahan lebih lama --- yaitu kesiagaan yang diberikan kopi.
\textbf{14.} $4000/50 = \qty{80}{s}$.
\textbf{15.} $600/5 = \qty{120}{s}$ tiap enzim tiap ribosom; 30 tiap
ribosom tiap jam, 600 tiap mRNA tiap jam.
\textbf{16.} $100\times 600 = 60\,000$ enzim tiap jam; $10^{5}$ setelah
kira-kira \qty{1.7}{h}.
\textbf{17.} Sepuluh sampai dua puluh menit bagi kortisol untuk masuk,
mengikat, dan mencapai gennya, sejam sebelum mRNA pertamanya banyak dan
tertranslasi, lalu dua jam penimbunan: kira-kira empat jam.
\textbf{18.} Sebuah riam mengubah keaktifan enzim yang sudah ada;
tugas kortisol adalah mengubah enzim mana yang ada, selama
berhari-hari, dan hanya transkripsi yang dapat mengerjakannya. Sebuah
ketakutan tidak dapat menunggu empat jam bagi enzim yang baru.
\textbf{19.} Reseptor yang lebih banyak berarti reseptor terisi yang
lebih banyak pada kepekatan adrenalin mana pun: bagian keterisiannya
tak berubah tetapi tanggapannya pada tiap kepekatannya lebih besar ---
kurvanya diskalakan naik, bukan digeser. Sebuah \omterm{def:b2:cell-signalling:kinds}{hormon} yang lambat
menetapkan bati \omterm{def:b2:cell-signalling:kinds}{hormon} yang cepat.
\textbf{20.} Glukosa masuk ke dalam sel $\beta$-nya lalu
dimetabolismekan; ATP-nya naik lalu menutup sebuah kanal kalium yang
peka ATP; membrannya berdepolarisasi; kanal kalsium bergerbang tegangan
membuka; kalsiumnya memicu eksositosis butiran insulinnya ---
metabolisme yang digandengkan kepada sekresi oleh sebuah langkah
kelistrikan.
\textbf{21.} Kedua reseptornya memulai dua lintasan yang bertemu pada
enzim yang sama: kinase A glukagonnya memfosforilasinya, riam insulinnya
mengaktifkan fosfatase yang mendefosforilasinya; keadaan enzimnya, dan
dengan demikian arah metabolisme glikogennya, mengikuti nisbah
keduanya.
\textbf{22.} Glukosanya ketika berpuasa tinggi (hatinya, yang tak
tertahan, terus membuatnya), glukosanya setelah santapan sangat tinggi
(tanpa penyerapan oleh otot dan lemaknya); tanpa insulin simpanan
lemaknya melepaskan asam lemak, yang diubah hatinya menjadi asam keton
--- ketoasidosis.
\textbf{23.} Otot yang bekerja menyerap glukosa tanpa insulin,
sehingga pada sebuah lari yang panjang glukosanya turun sementara
insulin yang disuntikkan, yang tak teratur, terus menekan hatinya:
hipoglikemia dan pingsan.
\textbf{24.} Gelung glukosanya: pengindranya sel $\beta$ dan sel
$\alpha$, dua \omterm{def:b2:cell-signalling:kinds}{hormon} efektor bagi kedua arahnya, bati yang tinggi yang
menahan arasnya sampai seperlima, beberapa menit untuk bekerja.
\omterm{def:b2:blood-pressure:baroreflex}{Barorefleks}: reseptor regangan, \omterm{def:b2:heart:anatomy}{jantung} dan pembuluhnya sebagai
efektor, bati kira-kira empat, sedetik untuk bekerja.
\textbf{25.} Keterisiannya $0.17$ setelah pelepasannya; batinya $5\times 10^{9}$ tiap detik
tiap reseptor yang terisi; keluaran hatinya \qty{5}{g/min}
sesungguhnya; kortisol memerlukan kira-kira empat jam.
\end{solution}
